\documentclass[conference]{IEEEtran}

\usepackage[utf8]{inputenc}
\usepackage{amsmath,amssymb,amsfonts}
\usepackage{graphicx}
\usepackage[hidelinks]{hyperref}
\usepackage{xurl}
\usepackage{cite}
\usepackage{balance}
\graphicspath{{./}{figures/}}
\DeclareGraphicsExtensions{.png,.pdf,.jpg}
\setlength{\textfloatsep}{3pt plus 1pt minus 1pt}
\setlength{\floatsep}{2.5pt plus 1pt minus 1pt}
\setlength{\intextsep}{3pt plus 1pt minus 1pt}
\setlength{\abovecaptionskip}{2pt}
\setlength{\belowcaptionskip}{0pt}
\setlength{\tabcolsep}{3pt}
\renewcommand{\arraystretch}{0.9}

\makeatletter
\newcommand{\linebreakand}{%
  \end{@IEEEauthorhalign}
  \hfill\mbox{}\par
  \mbox{}\hfill\begin{@IEEEauthorhalign}
}
\makeatother

\title{Neurodapt: \\
Cognitively-Aligned Reinforcement Learning for Memory-Optimized Note Rewriting}

\author{
\IEEEauthorblockN{Vinayak Mishra}
\IEEEauthorblockA{Department of CSE\\
AMC Engineering College\\
Bangalore, India\\
\texttt{vk2005917@gmail.com}}
\and
\IEEEauthorblockN{Vasisht Suresh}
\IEEEauthorblockA{Department of CSE\\
AMC Engineering College\\
Bangalore, India\\
\texttt{vasisht765@gmail.com}}
\and
\IEEEauthorblockN{Ullas Roy}
\IEEEauthorblockA{Department of CSE\\
AMC Engineering College\\
Bangalore, India\\
\texttt{ullasroy05@gmail.com}}
\linebreakand
\IEEEauthorblockN{Sushmitha N}
\IEEEauthorblockA{Department of CSE\\
AMC Engineering College\\
Bangalore, India\\
\texttt{nsushmitha2005@gmail.com}}
\and
\IEEEauthorblockN{R. Velvizhi Ramya}
\IEEEauthorblockA{Assistant Professor, Department of CSE\\
AMC Engineering College\\
Bangalore, India\\
\texttt{rvramya2021@gmail.com}}
}

\begin{document}

\maketitle

\begin{abstract}
Large Language Models (LLMs) are increasingly used to rewrite, summarize, and simplify text while preserving its meaning. Most of these systems, however, are optimized for fluency, semantic similarity, or task performance rather than for the information a reader is likely to retain over time. We present Neurodapt, a framework for long term memory retention oriented rewriting that combines Transformer-based clause boundary predictor, a memory-ranking module and a GRPO fine tuned LLM . The intended rewriting objective is limited to three components: improving the Memory Ranker score, preserving semantic meaning using a BERT-based measure, and enforcing a length constraint. We report component-level Clause Segmentation, memory-ranking and LLM rewriting results.
\end{abstract}

\begin{IEEEkeywords}
Large Language Models, Reinforcement Learning, Educational Technology, Memory Modeling, Cognitive Computing
\end{IEEEkeywords}

\section{Introduction}



LLMs have made rewriting, summarization, and educational content generation easier to automate. Yet the objectives used by many systems focus on linguistic quality or semantic preservation rather than on whether important information will be retained by the reader in the long term. A rewrite can therefore be faithful in meaning while still changing the cues that that can help the user's brain retain the information.

We propose \textbf{Neurodapt}, an approach to memory-oriented note rewriting. The framework first divides a note into clauses, estimates their relative memorability using a model trained on narrative recall data, and then uses those estimates to train a model to rewrite the text at a clause level. Its main contributions are:
\begin{enumerate}
\item A clause-boundary prediction module that uses a Transformer to estimate the probability of each token being a clause boundary.
\item A Transformer-based Memory Ranker trained from narrative retelling data to estimate relative clause memorability.
\item An LLM fine-tuned with GRPO to rewrite clauses toward higher memory-retention scores produced by the Memory Ranker.
\end{enumerate}

Together, these components form a pipeline for identifying clause boundaries, estimating their memory-retention scores, and optimizing clauses for improved retention.

\section{Related Work}

LLMs are now widely used for rewriting, summarization, and educational content generation \cite{zhang2023summit}. Much of this work emphasizes fluency, semantic similarity, or downstream task performance rather than the likelihood that particular information will be retained \cite{nakatani2022}. This distinction is important because a fluent rewrite can still alter the organization and cues that support later retrieval \cite{roit2023}.

Research on narrative memory offers a useful source of supervision for this problem. Georgiou et al. studied recall of meaningful narratives at scale and reported relationships between clause-level properties and human recall probabilities \cite{georgiou2025}. Cognitive accounts such as ACT-R describe retrieval in terms of activation, context, and interference \cite{bothell2021}, while spaced-repetition models focus on review history and forgetting \cite{settles2016}. Neurodapt differs in emphasis: instead of only modeling memory after information is presented, it uses memory-related signals to inform how the information is rewritten.

Embedding-based measures such as BERTScore are established tools for assessing semantic preservation \cite{zhang2019}. Neurodapt builds on this idea by combining semantic preservation with clause-level segmentation, narrative-based memorability estimation, and a reinforcement-learning formulation.

\section{Proposed Framework}

Given an input text, the framework first segments the text into individual clauses using the clause segmentation module. Each clause is then assigned a relative memory-retention score by the Memory Ranker, which estimates how strongly the clause is expected to contribute to the retention of the overall text.

Based on these scores, a target clause is selected for rewriting by the user. The target clause is provided to the rewriting LLM together with its surrounding clauses as contextual information. The rewriting LLM is fine-tuned using Group Relative Policy Optimization (GRPO), with the Memory Ranker providing the reward signal during training. This encourages the model to generate rewrites that achieve higher estimated memory-retention scores while remaining conditioned on the original context.

\begin{figure}[!h]
\centering
\includegraphics[width=\columnwidth]{Fig1.png}
\caption{Overall architecture of the Neurodapt framework.}
\label{fig:architecture}
\end{figure}

The complete framework therefore consists of three main components: clause segmentation, memory-retention ranking, and memory-guided clause rewriting.The implementation of these three components is described in the following sections.

\section{Clause Segmentation}

\subsection{Motivation and Problem Formulation}

We use clauses as the basic unit of memory analysis because they are more informative than individual tokens while remaining more specific than full sentences. A clause typically expresses an event, action, state, or proposition that can be treated as a distinct unit of information. A sentence may contain several such events, making sentence-level scoring too coarse for relative memorability. Clause segmentation therefore provides a practical middle ground between token-level and sentence-level representations \\cite{seyfried2024}.

\subsection{HCA Dataset and Preprocessing}

The clause boundary predictor is trained on the English HCA dataset (\texttt{eng.hca.amr}) from HCASeg \cite{yu2023}. The corpus is provided in CoNLL-U format, with clause boundaries marked by \texttt{BeginSeg=Yes} in the \texttt{MISC} field. We use the original training, development, and test partitions.

For a sentence
\begin{equation}
S=(w_1,w_2,\ldots,w_n),
\end{equation}
the clause annotations are converted into binary boundary labels
\begin{equation}
Y=(y_1,y_2,\ldots,y_n), \qquad y_i\in{0,1}.
\end{equation}
The annotated boundaries are aligned with BERT subword tokens through character offsets so that WordPiece tokenization does not shift the boundary labels.

\subsection{Contextual Representation and Boundary Prediction}

For each sentence, \texttt{bert-base-uncased} produces contextual token representations $h_i\in\mathbb{R}^{768}$. We mean-pool these representations to obtain a sentence-level representation and concatenate it with each token representation:

\begin{equation}
x_i=[h_i;h_S]\in\mathbb{R}^{1536}.
\end{equation}

Sequences are limited to 512 BERT tokens. The concatenated representations are projected to 384 dimensions and passed through a four-layer Transformer encoder with four attention heads, a feed-forward dimension of 1536, and dropout of 0.2:

\begin{equation}
z_i=W_px_i+b_p.
\end{equation}

A linear layer produces a boundary logit $l_i$, which is converted to a probability with the sigmoid function. We predict a boundary when

\begin{equation}
\hat{y}_i=
\begin{cases}
1, & p_i>\tau,\\
0, & p_i\leq\tau.
\end{cases}
\end{equation}

A threshold of $\tau=0.39$ is selected by validation-based calibration.

\subsection{Training Objective}

We use binary cross-entropy for boundary prediction, excluding padding positions from the loss. The model is trained with AdamW using a learning rate of $10^{-4}$, batch size 4, and up to 30 epochs. We retain the checkpoint with the lowest development loss.

Validation threshold calibration is applied after training to select the decision threshold used at test time.

\subsection{Clause Reconstruction and Post-Processing}

During inference, the predicted boundaries are combined with the BERT tokens to reconstruct contiguous clauses. Clauses containing tokens with \texttt{\#\#} are removed. The remaining tokens are grouped until a predicted boundary is reached, producing the ordered clause sequence

\begin{equation}
C=(c_1,c_2,\ldots,c_k).
\end{equation}

A lightweight post-processing step is then applied to make the reconstructed clauses linguistically usable. We remove WordPiece markers, merge contractions, and normalize spacing around punctuation, slashes, and hyphens. Punctuation-only segments are attached to adjacent clauses, while single-word segments are merged with neighboring clauses.

Finally, clauses containing fewer than four whitespace-delimited tokens are merged iteratively with an adjacent clause. A short first or last clause is merged with its only neighbor. For an intermediate short clause, we merge it with the smaller neighboring clause unless the left neighbor ends with a period, in which case it is merged with the right neighbor. The resulting clause sequence is then passed to the memory-ranking module.

\subsection{Architecture}

The complete training and inference architecture of the clause segmentation module is shown in Fig.~\ref{fig:clause_architecture}.

\begin{figure}[!t]
\centering
\includegraphics[width=\columnwidth,height=1.90in,keepaspectratio]{Fig2.png}
\caption{Training and inference architecture of the clause segmentation module.}
\label{fig:clause_architecture}
\end{figure}


\section{Memory Ranking Module}

After segmentation, the next module estimates the relative memorability of individual clauses using narrative recall data.

\subsection{Narrative Recall Data}

The memory-ranking model is trained using the Breithaupt narrative dataset \cite{he2023}, which contains an original story and three subsequent human retellings. We pass the original story and each retelling through the trained clause segmenter to obtain ordered clause sequences.

For an original story, let

\begin{equation}
C^{(o)}=(c^{(o)}_1,c^{(o)}_2,\ldots,c^{(o)}_N)
\end{equation}

denote its original clauses and let

\begin{equation}
C^{(g)}=(c^{(g)}_1,c^{(g)}_2,\ldots,c^{(g)}_{M_g})
\end{equation}

denote the clauses in retelling $g$, where $g\in\{1,2,3\}$.

\subsection{Clause-Level Recall Matching}

Each clause is represented by a normalized 768-dimensional \texttt{all-mpnet-base-v2} embedding. For an original clause $c^{(o)}_i$, we compute cosine similarity against the clauses in each retelling, and the strongest match becomes the retelling-level score:

\begin{equation}
r_i^{(g)}
=
\max_{j\in C^{(g)}}e_i^\mathsf{T}e_j^{(g)},
\end{equation}

where $e_i$ and $e_j^{(g)}$ are normalized clause embeddings.

The three retelling scores are summed:

\begin{equation}
r_i =
\sum_{g=1}^{3}r_i^{(g)}.
\end{equation}

A clause therefore receives a higher score when its content is reflected in the human retellings. The implementation also computes normalized source-clause centrality for inspection, but does not use it in the training target. The centrality score for a story with more than one clause is
\begin{equation}
a_i=\frac{1}{N-1}\sum_{j\ne i}e_i^\mathsf{T}e_j.
\end{equation}
For a single-clause story, centrality is set to one. The retelling-only target is min--max normalized within each story, as described next; if all retelling scores are equal, the target is set to 0.5 for every clause. No centrality/retelling mixture or clipping operation is used by the current data-preparation script.

\subsection{Within-Story Normalization}

The goal is to compare clauses within each story rather than compare absolute scores across different stories. The retelling scores are min--max normalized within each story. For a story with target range $[r_{\min},r_{\max}]$, the normalized score is

\begin{equation}
m_i =
\frac{r_i-r_{\min}}{r_{\max}-r_{\min}},
\end{equation}

If a story contains only one clause, or all clauses have the same score, every clause is assigned the neutral target value $0.5$.

This normalization preserves the within-story ordering while reducing the influence of differences in the absolute embedding-similarity scale.

\subsection{Transformer-Based Memory Ranker}

The normalized retelling scores are used as targets for a Transformer-based Memory Ranker. Each 768-dimensional clause embedding $x_i$ is first projected into a 256-dimensional hidden space:

\begin{equation}
z_i=W_px_i+b_p.
\end{equation}

A learned positional embedding captures the location of each clause in sequences of up to 512 clauses. The resulting sequence is processed by two Transformer encoder layers with four attention heads, a feed-forward dimension of 1024, GELU activation, and dropout of 0.1. A two-layer MLP with ReLU and dropout 0.1 produces one unbounded logit per clause; sigmoid is applied for probability-scale predictions.

\subsection{Ranking Objective}

Because the primary goal is relative ranking within a story, the training objective combines binary cross-entropy with a margin-ranking loss:

\begin{equation}
\mathcal{L}=0.5\mathcal{L}_{\mathrm{BCE}}+0.5\mathcal{L}_{\mathrm{rank}},
\end{equation}

where the pointwise term is binary cross-entropy with logits, and, for unequal target scores, the pairwise term uses margin 0.2:

\begin{equation}
\ell_{ij}=\max\left(0,0.2-\operatorname{sign}(m_i-m_j)(\hat{s}_i-\hat{s}_j)\right).
\end{equation}

The pairwise losses are averaged within each story and then across the batch. Variable-length stories are padded and passed with a key-padding mask; tied target pairs are excluded from the ranking loss. Sigmoid is applied to raw model logits for probability-scale inference, not in the pairwise loss.

\subsection{Training and Evaluation}

The Memory Ranker is trained with AdamW at a learning rate of $10^{-4}$, batch size 8, and a maximum of 50 epochs. A ReduceLROnPlateau scheduler halves the learning rate after two non-improving validation epochs, down to a minimum of $10^{-6}$. Stories are shuffled deterministically using seed 42 and split by story into 70\% training, 10\% validation, and the remaining test set (integer truncation determines the exact counts). Early stopping monitors validation loss with patience 10, and the lowest-loss checkpoint is saved.

For evaluation, we use the independent human-memory dataset reported by Georgiou et al. \cite{georgiou2025}, which provides human-derived clause recall probabilities. Because the benchmark is small, it is used only for evaluation. We compare predicted memorability with human recall using MAE, RMSE, Spearman correlation, Kendall's $\tau$, pairwise ranking accuracy, NDCG@1, NDCG@3, NDCG@5, and top-5 overlap. Metrics are computed separately for each story and then averaged across stories. The comparison includes random ranking, clause position, and MPNet embedding similarity computed from the same normalized representation.

\subsection{Target-Ratio Ablation}

The current data-preparation script trains on the retelling-only target and does not implement the earlier centrality--retelling ratio sweep. The sweep values and associated figures are therefore omitted from the current experimental claims. Centrality is computed and retained as a diagnostic field, but a controlled ablation would require regenerating targets, retraining each configuration, and evaluating each resulting checkpoint with the same protocol.

\subsection{Memory Ranking Pipeline}

Fig.~\ref{fig:memory_ranker} summarizes the training and inference pipeline of the Memory Ranking Module.

\begin{figure}[!t]
\centering
\includegraphics[width=\columnwidth,height=1.90in,keepaspectratio]{Fig3.png}
\caption{Training and inference pipeline of the Memory Ranking Module.}
\label{fig:memory_ranker}
\end{figure}

\section{Reinforcement Learning Pipeline}

The intended RL reward is restricted to three components: (1) a Memory Ranker score for memorability, (2) a BERT-based semantic-preservation score comparing the source target clause with its rewrite, and (3) a length constraint. There is no semantic-centrality reward or separate information-fidelity reward. The current trainer provides a GRPO experiment for rewriting one target clause in its surrounding clause context, but currently implements only the Memory Ranker component and the length constraint; the BERT-based semantic-preservation reward is planned, not yet implemented. The training data are produced from 50-word windows with stride 50. The existing clause segmenter and Memory Ranker score each window; interior clauses are selected in center-out order as target examples. The prompt instructs the model to preserve the target meaning, avoid adding information or borrowing from other clauses, produce one sentence, and stay within 150\% of the original target's word count. Up to 1,500 valid examples are selected using seed 42 by default; the preprocessing pipeline does not define train/validation/test partitions.

The policy is the local Qwen2.5-0.5B-Instruct model, trained for one epoch with TRL's GRPO trainer. The run uses 4-bit NF4 quantization with double quantization, FP16 computation, and LoRA adapters (rank 8, alpha 16, dropout 0.05) on the query, key, value, and output projections. The per-device batch size is 2, gradient accumulation is 8, four completions are sampled per prompt, the learning rate is $5\times10^{-6}$, and the maximum prompt and completion lengths are 256 and 48 tokens. The saved run contains 5,000 training steps and reaches epoch 1.

For each candidate, the reward scorer embeds the original target and rewrite with normalized MPNet embeddings and scores them together as a two-clause sequence. Let $\Delta=s_{\mathrm{rewrite}}-s_{\mathrm{original}}$, where each $s$ is the Memory Ranker's sigmoid output, and let $L_o$ and $L_w$ be the original and rewritten word counts. The implemented reward is
\begin{equation}
R=\begin{cases}
-2.5, & L_o=0\text{ or }L_w=0,\\
\Delta, & 2L_w\leq 3L_o,\\
-1.5-\left(\frac{L_w}{1.5L_o}-1\right)+0.25\Delta, & 2L_w>3L_o.
\end{cases}
\end{equation}
Thus valid pairwise score differences lie in $[-1,1]$, while completions exceeding the 150\% word-count cap receive a rejection penalty plus an excess-length penalty and a small ranker-score contribution. The BERT-based semantic-preservation score will be added to this implemented ranker-and-length reward; it is not reflected in the checkpoint's training objective. In the intended three-component objective, the final reward will combine these ranker, BERT-semantic, and length-constraint signals, with no additional centrality or information-fidelity term.

The checkpoint and trainer log establish that optimization ran, but are not evidence of improved retention, semantic preservation, or overall rewrite quality. The available comparison utility supports qualitative base-model versus adapter generations with ranker scoring; no held-out quantitative evaluation of the GRPO outputs is currently provided. Integrating the BERT-based reward and evaluating the complete three-component objective remain future work.

\section{Experiments}

\subsection{Experimental Setup}

The experiments report held-out clause segmentation results and Memory Ranker performance on 13 independent stories with human clause-recall probabilities, using MAE, RMSE, Spearman correlation, Kendall's $\tau$, pairwise accuracy, NDCG@1/3/5, and top-5 overlap. The Memory Ranker is compared with random ranking, clause position, and MPNet embedding similarity. GRPO training has also been run and produced a saved adapter checkpoint, but the present study does not report a controlled quantitative evaluation of its generated rewrites or reader retention.

\subsection{GRPO Training Run}

The GRPO experiment was designed as a controlled first test of whether a memorability-guided reward produces observable changes in a small instruction-tuned policy. The raw fine-tuning data were converted to CSV, and rows containing fewer than 50 whitespace-delimited words were removed. Each remaining text was divided into non-overlapping 50-word windows using a stride of 50. The clause segmenter and Memory Ranker were then applied to each window. Windows with fewer than three clauses were discarded, and the remaining interior clauses were selected in center-out order as target clauses. Each target produced one training example containing the ordered clause context, the target marker, the original target text, and its Memory Ranker score. This procedure generated 198,800 validatable examples in the stored preprocessing artifact.

The trainer used a deterministic sample of 20,000 examples from these rows, selected with seed 42. This cap was chosen because the purpose of this run was to establish whether the reward and prompting scheme could produce measurable behavioral changes, rather than to maximize corpus coverage or claim dataset-size saturation. Using more windows was therefore not necessary for this initial diagnostic experiment and would have increased compute and training time without changing the question being tested. The 20,000 examples were used for the single training run; the preprocessing and trainer do not define a held-out validation or test partition, so the run cannot support generalization claims.

Each prompt asked the policy to rewrite only the marked target clause, preserve its meaning, avoid information from neighboring clauses, produce one sentence, and remain within 150\% of the original target's word count. The local Qwen2.5-0.5B-Instruct model was loaded with 4-bit NF4 double-quantized weights and FP16 computation. LoRA was applied to the query, key, value, and output projections with rank 8, alpha 16, and dropout 0.05. TRL's GRPO configuration used a per-device batch size of 2, gradient accumulation of 8, four sampled completions per prompt, one epoch, a learning rate of $5\times10^{-6}$, gradient checkpointing, and a maximum completion length of 48 tokens. The maximum prompt length was 256 tokens. No external experiment tracker was used; metrics were written to the trainer log.

The implemented reward uses the Memory Ranker score change for the target clause and a length constraint. If $s_o$ and $s_w$ are the ranker scores for the original and rewritten target in the same clause context, the score contribution is $\Delta=s_w-s_o$. Empty outputs receive a fixed penalty of $-2.5$. Rewrites within the 150\% word-count limit receive $\Delta$; longer rewrites receive a rejection penalty, an excess-length penalty, and a reduced $0.25\Delta$ contribution. The BERT-based semantic-preservation term described in the intended objective was not present in this run. Consequently, the run tests ranker-guided rewriting under a length constraint, not the complete three-component objective.

The logged run completed one epoch in approximately $27{,}250$ seconds, corresponding to 5,000 training steps. The final trainer summary reports a training loss of 0.06289, while the last logged interval reports a reward mean of approximately $-0.184$, reward standard deviation 0.265, mean completion length 10.88 tokens, and clipped-completion ratio 0.0125. These values describe the optimization trajectory and sampling behavior only. In particular, the negative reward scale is expected under the implemented empty/overlength penalties and does not by itself imply failed learning. The run does not include held-out rewrite-quality, semantic-fidelity, or human-retention measurements.

\begin{figure}[!t]
\centering
\includegraphics[width=\columnwidth]{fine_tuning/outputs/Qwen2.5-0.5-Instruct-GRPO/plots/loss_reward.png}
\caption{GRPO loss and reward recorded during the 20,000-example training run.}
\label{fig:grpo_loss_reward}
\end{figure}

\begin{figure}[!t]
\centering
\includegraphics[width=\columnwidth]{fine_tuning/outputs/Qwen2.5-0.5-Instruct-GRPO/plots/reward_statistics.png}
\includegraphics[width=\columnwidth]{fine_tuning/outputs/Qwen2.5-0.5-Instruct-GRPO/plots/completion_lengths.png}
\caption{Reward variability and generated completion lengths during GRPO training.}
\label{fig:grpo_reward_lengths}
\end{figure}

\begin{figure}[!t]
\centering
\includegraphics[width=\columnwidth]{fine_tuning/outputs/Qwen2.5-0.5-Instruct-GRPO/plots/learning_rate.png}
\includegraphics[width=\columnwidth]{fine_tuning/outputs/Qwen2.5-0.5-Instruct-GRPO/plots/entropy_gradient_norm.png}
\includegraphics[width=\columnwidth]{fine_tuning/outputs/Qwen2.5-0.5-Instruct-GRPO/plots/clipped_ratio.png}
\caption{Learning-rate decay, entropy and gradient-norm behavior, and completion clipping during the run.}
\label{fig:grpo_training_diagnostics}
\end{figure}

\subsection{Clause Segmentation Results}

Clause segmentation is evaluated using precision, recall, accuracy, and F1. The fixed-learning-rate configuration reaches a test F1 of 0.8153, while adding a learning-rate scheduler increases it to 0.8252.

\begin{table}[t]
\centering
\caption{Clause segmentation performance under different configurations.}
\label{tab:clause_results}
\begin{tabular}{lcccc}
\hline
Configuration & Precision & Recall & Accuracy & F1 \\
\hline
Fixed LR, $\tau=0.50$ & 0.7808 & 0.8530 & 0.9737 & 0.8153 \\
LR Scheduler, $\tau=0.50$ & 0.7900 & 0.8637 & 0.9751 & 0.8252 \\
LR Scheduler, $\tau=0.39$ & \textbf{0.8481} & 0.8099 & \textbf{0.9772} & \textbf{0.8285} \\
\hline
\end{tabular}
\end{table}

\subsubsection{Learning Rate Scheduling}

With the other settings unchanged, adding the learning-rate scheduler increases test F1 from 0.8153 to 0.8252. This improvement suggests better generalization than the fixed-learning-rate setting.

\subsubsection{Decision Threshold Calibration}

We tune the decision threshold on the validation set over 0.10--0.90 in steps of 0.01. The best validation F1 occurs at $\tau=0.39$, giving F1 = 0.8487, precision = 0.8632, recall = 0.8347, and accuracy = 0.9804.

\begin{figure}[!t]
\centering
\includegraphics[width=\columnwidth,height=1.80in,keepaspectratio]{Fig_Threshold_F1.png}
\caption{Validation F1 across decision thresholds; the selected threshold is $\tau=0.39$.}
\label{fig:threshold_f1}
\end{figure}

We then fix this threshold and apply it to the held-out test set. The resulting precision is 0.8481, recall 0.8099, accuracy 0.9772, and F1 0.8285. Relative to the default threshold of 0.50, calibration improves precision and gives a modest F1 gain.

\subsection{Memory Ranking Results}

Table~\ref{tab:memory_results} reports the exact evaluation output. The Memory Ranker improves on all listed baselines for MAE, RMSE, Spearman correlation, Kendall's $\tau$, pairwise accuracy, and all NDCG values. Its top-5 overlap of 0.2308 ties the random baseline and exceeds the position and MPNet baselines.

The benchmark contains only 13 stories, so these results support a useful ranking signal but do not establish a complete model of human memory.

\begin{table}[t]
\centering
\caption{Memory ranking performance on the independent human-memory evaluation set. Lower is better for MAE and RMSE; higher is better for all other metrics.}
\label{tab:memory_results}
\resizebox{\columnwidth}{!}{%
\begin{tabular}{lccccccccc}
\hline
Model & MAE & RMSE & Spearman & Kendall $\tau$ &
Pairwise Acc. & NDCG@1 & NDCG@3 & NDCG@5 & Top-5 \\
\hline
Memory Ranker & \textbf{0.2414} & \textbf{0.2923} & \textbf{0.2797} & \textbf{0.2097} &
\textbf{0.6049} & \textbf{0.4999} & \textbf{0.5021} & \textbf{0.5232} & \textbf{0.2308} \\
Random & 0.3486 & 0.4202 & -0.0497 & -0.0436 &
0.4782 & 0.3728 & 0.4436 & 0.4690 & \textbf{0.2308} \\
Position & 0.3621 & 0.4332 & -0.1096 & -0.0801 &
0.4599 & 0.4032 & 0.4118 & 0.4378 & 0.1077 \\
MPNet Embedding & 0.3518 & 0.4063 & 0.1282 & 0.0822 &
0.5411 & 0.4883 & 0.4482 & 0.4735 & 0.1385 \\
\hline
\end{tabular}}
\end{table}

Figures~\ref{fig:ranking_metrics}--\ref{fig:error_metrics} summarize the ranking and prediction-error comparisons.

\begin{figure}[!t]
\centering
\includegraphics[width=\columnwidth,height=1.85in,keepaspectratio]{Fig4_RankingMetrics.png}
\caption{Comparison of ranking performance across the Memory Ranker and baseline methods.}
\label{fig:ranking_metrics}
\end{figure}

\begin{figure}[!t]
\centering
\includegraphics[width=\columnwidth,height=1.85in,keepaspectratio]{Fig5_NDCG.png}
\caption{NDCG performance at different ranking depths.}
\label{fig:ndcg_metrics}
\end{figure}

\begin{figure}[!t]
\centering
\includegraphics[width=\columnwidth,height=1.70in,keepaspectratio]{Fig6_ErrorMetrics.png}
\caption{Prediction error relative to human clause-recall probabilities.}
\label{fig:error_metrics}
\end{figure}


\section{Discussion}

Taken together, the component-level experiments support the feasibility of the proposed pipeline. Clause segmentation reaches an F1 of 0.8285 on the held-out HCA test set, and the Memory Ranker improves on the listed baselines on several ranking measures, including pairwise accuracy and the NDCG metrics.

The benchmark also highlights an important limitation: it contains only 13 stories, and the training targets are indirect proxies formed from retelling similarity. The current Memory Ranker is therefore better interpreted as a useful ranking signal, not as a complete model of human memory.

A second limitation concerns how the training targets are constructed. They are inferred by matching original clauses to human retellings rather than obtained from direct clause-level retention labels. The independent human-memory benchmark provides an external check, but its small size means that broader claims about human memory should be made cautiously. The current maintained pipeline uses retelling-only targets; it does not support conclusions about the comparative contribution of semantic centrality because the earlier ratio sweep is not reproduced by the current target-generation script.

The current GRPO implementation is a limited experiment rather than a complete evaluation of the Neurodapt hypothesis: its reward optimizes Memory Ranker score change for one target clause and penalizes empty or overlength completions. The intended BERT-based semantic-preservation component has not yet been integrated. The saved adapter confirms that training ran, but the work does not yet report held-out measurements of semantic preservation, rewrite quality, or human retention. These evaluations should be conducted after integrating the full intended three-component reward.

\section{Conclusion}

This paper introduced Neurodapt, a modular framework that combines clause-level segmentation, narrative-based memorability estimation, and GRPO-based target-clause rewriting. The segmentation model reaches an F1 of 0.8285, while the independent human-memory benchmark shows that the retelling-trained Memory Ranker outperforms listed baselines on several metrics. The ranker is trained on within-story normalized retelling-match scores; the current implementation does not use a centrality mixture, and the earlier ratio sweep is not reproduced by the maintained pipeline. The current GRPO trainer uses a ranker-score-difference reward with a length constraint; the intended BERT-based semantic-preservation term has not yet been integrated. The effect on semantic preservation, generated-text quality, and human retention has not yet been measured. The small benchmark and absence of a controlled rewrite evaluation limit conclusions about broader human-memory alignment and end-to-end benefit.

These findings motivate testing human-aligned memorability signals as an objective for LLM-based rewriting. Future work should integrate the BERT-based semantic-preservation reward and evaluate the resulting three-component objective (Memory Ranker, BERT semantic preservation, and length constraint) against controlled baselines using rewrite quality and direct human-retention measures.

\balance
\small
\begin{thebibliography}{99}

\bibitem{georgiou2025}
A. Georgiou, T. Can, M. Katkov, and M. Tsodyks,
``Large-scale study of human memory for meaningful narratives,''
\emph{Learning \& Memory}, vol. 32, no. 2, Art. no. a054043, 2025,
doi: 10.1101/lm.054043.124.

\bibitem{seyfried2024}
C. Seyfried, C. Reed, and Y. Kamide,
``Defining argumentative discourse units as clauses: psycholinguistic evidence,''
in \emph{Computational Models of Argument: Proceedings of COMMA 2024},
SAGE Publications, 2024, pp. 277--288,
doi: 10.3233/FAIA240328.

\bibitem{bothell2021}
D. Bothell,
\emph{ACT-R 7.21+ Reference Manual},
Carnegie Mellon University,
Pittsburgh, PA, USA,
2021.

\bibitem{clark1991}
J. M. Clark and A. Paivio,
``Dual Coding Theory and Education,''
\emph{Educational Psychology Review},
vol. 3, no. 3, pp. 149--210, 1991.

\bibitem{he2023}
T. He, F. Breithaupt, S. Kübler, and T. T. Hills,
``Quantifying the retention of emotions across story retellings,''
\emph{Scientific Reports},
vol. 13, Art. no. 2448, 2023.

\bibitem{anderson2004}
J. R. Anderson et al.,
``An Integrated Theory of the Mind,''
\emph{Psychological Review},
vol. 111, no. 4, pp. 1036--1060, 2004.

\bibitem{labov1966}
W. Labov and J. Waletzky,
``Narrative Analysis: Oral Versions of Personal Experience,''
in \emph{Essays on the Verbal and Visual Arts},
Seattle, WA, USA: Univ. of Washington Press, 1966.

\bibitem{shao2024}
Z. Shao et al.,
``DeepSeekMath: Pushing the Limits of Mathematical Reasoning in Open Language Models,''
\emph{arXiv preprint arXiv:2402.03300},
2024.

\bibitem{zhang2019}
T. Zhang, V. Kishore, F. Wu, K. Q. Weinberger, and Y. Artzi,
``BERTScore: Evaluating Text Generation with BERT,''
in \emph{Proc. ICLR},
2019.

\bibitem{brysbaert2014}
M. Brysbaert, A. B. Warriner, and V. Kuperman,
``Concreteness Ratings for 40 Thousand Generally Known English Word Lemmas,''
\emph{Behavior Research Methods},
vol. 46, no. 3, pp. 904--911, 2014.

\bibitem{wartena2021}
C. Wartena,
``On the Geometry of Concreteness,''
in \emph{Proc. ACL},
2021.

\bibitem{settles2016}
B. Settles and B. Meeder,
``A Trainable Spaced Repetition Model for Language Learning,''
in \emph{Proc. ACL},
2016, pp. 1848--1858.

\bibitem{vaswani2017}
A. Vaswani, N. Shazeer, N. Parmar, J. Uszkoreit,
L. Jones, A. N. Gomez, {\L}. Kaiser, and I. Polosukhin,
``Attention Is All You Need,''
in \emph{Advances in Neural Information Processing Systems (NeurIPS)},
vol. 30, 2017.

\bibitem{kintsch1978}
W. Kintsch and T. A. van Dijk,
``Toward a Model of Text Comprehension and Production,''
\emph{Psychological Review},
vol. 85, no. 5, pp. 363--394, 1978.

\bibitem{zhang2023summit}
H. Zhang, X. Liu, and J. Zhang,
``Summit: Iterative Text Summarization via ChatGPT,''
in \emph{Findings of ACL: EMNLP},
2023, pp. 10644--10657.

\bibitem{nakatani2022}
Y. Nakatani and T. Kajiwara,
``Comparing BERT-based Reward Functions for Deep Reinforcement Learning in Machine Translation,''
in \emph{Proc. WAT},
2022, pp. 37--44.

\bibitem{zhang2025qapyramid}
S. Zhang et al.,
``QAPyramid: Fine-grained Evaluation of Content Selection for Text Summarization,''
in \emph{Second Conference on Language Modeling},
2025.

\bibitem{roit2023}
P. Roit et al.,
``Factually Consistent Summarization via Reinforcement Learning with Textual Entailment Feedback,''
in \emph{Proc. ACL},
2023.

\bibitem{qwen2024}
Qwen Team,
``Qwen2.5 Technical Report,''
\emph{arXiv preprint arXiv:2412.15115},
2024.

\bibitem{yu2023}
Y. Fan, B. Li, Y. Sataer, M. Gao, C. Shi, S. Cao, and Z. Gao,
``Hierarchical Clause Annotation: Building a Clause-Level Corpus for Semantic Parsing with Complex Sentences,''
\emph{Applied Sciences}, vol. 13, no. 16, Art. no. 9412, 2023,
doi: 10.3390/app13169412.

\bibitem{openai2024embeddings}
OpenAI,
``New Embedding Models and API Updates,''
2024.
[Online]. Available:
\url{https://openai.com/index/new-embedding-models-and-api-updates/}
\end{thebibliography}

\end{document}
